% !PRINTED 367
\weberpage{367}
\S.\,98.

\gesperrt{Die Aronhold'schen Siebener-Systeme.}

Von besonderer Wichtigkeit sind die azygetischen Systeme von sieben Doppeltangenten, die zuerst von \gesperrt{Aronhold} betrachtet sind, und die daher \gesperrt{Aronhold'sche Siebener-Systeme} heissen. Wir nennen sie auch \gesperrt{vollständige Siebener-Systeme} oder kurz \gesperrt{vollständige Systeme}.

Dass es solche Systeme giebt, zeigen die Zusammenstellungen des vorigen Paragraphen. Denn wenn

\weq{1}{\begin{aligned}
&x_1 y_1,\ x_2 y_2,\ x_3 y_3,\ x_4 y_4,\ x_5 y_5,\ x_6 y_6\\
&x_1 z_1,\ x_2 z_2,\ x_3 z_3,\ x_4 z_4,\ x_5 z_5,\ x_6 z_6
\end{aligned}}

ein azygetisches Complexpaar ist, so ist

\[ x_1,\ x_2,\ x_3,\ x_4,\ x_5,\ y_6,\ z_6 \]

ein vollständiges Siebener-System (weil in dem Complexe $y_6 z_6$ keines der $x$ vorkommt). Es wird sich zeigen, dass keine azygetischen Systeme von mehr als sieben Elementen bestehen. Es gilt zunächst der Satz:

7. \gesperrt{Irgend sechs Elemente eines vollständigen Systemes kommen in einem und nur in einem Complexe vor.}

Wir beweisen zunächst den zweiten Theil der Behauptung, d.\,h. wir beweisen, dass, wenn

\weq{2}{x_1,\ x_2,\ x_3,\ x_4,\ x_5,\ x_6,\ x_7}

ein vollständiges System ist, $x_1, x_2, x_3, x_4, x_5, x_6$ nicht in zwei verschiedenen Complexen vorkommen können. Nehmen wir an, es sei dies möglich, so müssen die beiden Complexe ein azygetisches Paar wie (1) bilden. In dem syzygetischen Tripel

\[\begin{aligned}
&y_1 z_1,\ y_2 z_2,\ y_3 z_3,\ y_4 z_4,\ y_5 z_5,\ y_6 z_6\\
&y_1 y_2,\ z_1 z_2,\ x_1 x_2\\
&y_1 z_2,\ y_2 z_1
\end{aligned}\]

müssen die Doppeltangenten $x_3, x_4, x_5, x_6, x_7$ vorkommen, und da sie weder im ersten noch im zweiten dieser Complexe vorkommen, so müssen sie im dritten vorkommen. Da in diesem Complexe aber nur noch vier Paare übrig sind, so müssen minde-
%CONT>>
